\section{Python Fundamentals for Linear Programming}

Before diving into optimization modeling with PuLP, we need to establish the fundamental Python concepts that will be essential for your linear programming work. This section covers the core data structures, control flow, and packages you'll use regularly.

\subsection{Essential Data Structures}

\subsubsection{Lists}
Lists are ordered collections that can hold different types of data. They're particularly useful for storing sets of decision variables, constraints, or data points.

\begin{lstlisting}[language=Python]
# Creating lists for LP problems
decision_vars = ['x1', 'x2', 'x3']
supply = [100, 150, 200]
demand = [80, 120, 170]

# List operations useful in LP
costs = [3, 5, 2, 4]
costs.append(6)  # Add new cost
total_supply = sum(supply)  # Sum all elements
max_cost = max(costs)  # Maximum element

# Accessing elements (0-indexed)
first_var = decision_vars[0]  # 'x1'
last_supply = supply[-1]     # 200
middle_costs = costs[1:3]    # [5, 2]
\end{lstlisting}

\subsubsection{Dictionaries}
Dictionaries store key-value pairs and are excellent for representing parameters, variable coefficients, and solution values in optimization problems.

\begin{lstlisting}[language=Python]
# Transportation problem data
supply = {'Plant1': 100, 'Plant2': 150, 'Plant3': 200}
demand = {'Market1': 80, 'Market2': 120, 'Market3': 170}

# Cost matrix as nested dictionary
costs = {
    ('Plant1', 'Market1'): 3,
    ('Plant1', 'Market2'): 5,
    ('Plant2', 'Market1'): 4,
    ('Plant2', 'Market2'): 2
}

# Alternative representation
costs_by_plant = {
    'Plant1': {'Market1': 3, 'Market2': 5},
    'Plant2': {'Market1': 4, 'Market2': 2}
}

# Accessing values
plant1_supply = supply['Plant1']
cost_p1_m1 = costs[('Plant1', 'Market1')]
cost_alt = costs_by_plant['Plant1']['Market1']

# Useful dictionary operations
plants = list(supply.keys())
total_supply = sum(supply.values())
\end{lstlisting}

\subsubsection{Strings}
Strings are crucial for naming variables, creating output, and file handling in optimization problems.

\begin{lstlisting}[language=Python]
# Variable naming in optimization
base_name = "transport"
plant = "Plant1"
market = "Market2"

# String formatting for variable names
var_name = f"{base_name}_{plant}_{market}"  # "transport_Plant1_Market2"
var_name_alt = "{}_{}_{}".format(base_name, plant, market)

# Multiple variable names
plants = ['Plant1', 'Plant2', 'Plant3']
var_names = [f"supply_{plant}" for plant in plants]
# Result: ['supply_Plant1', 'supply_Plant2', 'supply_Plant3']

# String operations for output formatting
solution_status = "Optimal"
objective_value = 1250.75
report = f"Status: {solution_status}, Objective: ${objective_value:.2f}"
\end{lstlisting}

\subsection{NumPy Arrays}
NumPy arrays provide efficient numerical computing capabilities essential for linear algebra operations in optimization.

\begin{lstlisting}[language=Python]
import numpy as np

# Creating arrays from lists
costs_list = [3, 5, 2, 4]
costs_array = np.array(costs_list)

# 2D arrays for matrices (constraint coefficients, cost matrices)
cost_matrix = np.array([
    [3, 5, 2],
    [4, 2, 6],
    [1, 3, 4]
])

# Array operations useful in LP
constraint_matrix = np.array([
    [1, 1, 0],  # First constraint
    [0, 1, 1],  # Second constraint
    [1, 0, 1]   # Third constraint
])

rhs = np.array([100, 150, 200])  # Right-hand side values

# Matrix operations
solution = np.array([50, 30, 70])
objective_value = np.dot(costs_array[:3], solution)  # Linear combination

# Useful array methods
print(f"Cost matrix shape: {cost_matrix.shape}")
print(f"Minimum cost: {cost_matrix.min()}")
print(f"Row sums: {constraint_matrix.sum(axis=1)}")
\end{lstlisting}

\subsection{Pandas DataFrames}
DataFrames excel at handling tabular data, making them perfect for LP problem data management and solution analysis.

\begin{lstlisting}[language=Python]
import pandas as pd

# Creating DataFrames for optimization data
plants_data = pd.DataFrame({
    'Plant': ['Plant1', 'Plant2', 'Plant3'],
    'Supply': [100, 150, 200],
    'Fixed_Cost': [1000, 1200, 800]
})

markets_data = pd.DataFrame({
    'Market': ['Market1', 'Market2', 'Market3'],
    'Demand': [80, 120, 170],
    'Price': [15, 18, 12]
})

# Transportation costs as DataFrame
import itertools
routes = list(itertools.product(plants_data['Plant'], markets_data['Market']))
transportation_df = pd.DataFrame({
    'From': [route[0] for route in routes],
    'To': [route[1] for route in routes],
    'Cost': [3, 5, 2, 4, 2, 6, 1, 3, 4]  # Example costs
})

# Useful DataFrame operations for LP
print(plants_data.head())
total_supply = plants_data['Supply'].sum()
high_cost_routes = transportation_df[transportation_df['Cost'] > 3]

# Pivoting for matrix representation
cost_matrix_df = transportation_df.pivot(index='From', columns='To', values='Cost')
print(cost_matrix_df)
\end{lstlisting}

\subsection{Control Structures}

\subsubsection{For Loops}
For loops are essential for iterating over decision variables, constraints, and data in optimization problems.

\begin{lstlisting}[language=Python]
# Basic for loop over lists
plants = ['Plant1', 'Plant2', 'Plant3']
supply = [100, 150, 200]

# Loop with index and value
for i, plant in enumerate(plants):
    print(f"{plant} has supply capacity: {supply[i]}")

# Loop over dictionary items
supply_dict = {'Plant1': 100, 'Plant2': 150, 'Plant3': 200}
for plant, capacity in supply_dict.items():
    print(f"{plant}: {capacity}")

# Nested loops for constraint generation
for i, plant in enumerate(plants):
    for j, market in enumerate(['Market1', 'Market2']):
        var_name = f"x_{i}_{j}"
        print(f"Variable {var_name}: {plant} to {market}")

# Loop with conditional logic
total_capacity = 0
for plant, capacity in supply_dict.items():
    if capacity > 120:
        print(f"{plant} is a large facility")
        total_capacity += capacity
\end{lstlisting}

\subsubsection{List Comprehensions}
List comprehensions provide a concise way to create lists, particularly useful for generating variable names, constraint expressions, and processing optimization results.

\begin{lstlisting}[language=Python]
# Basic list comprehension
plants = ['Plant1', 'Plant2', 'Plant3']
markets = ['Market1', 'Market2', 'Market3']

# Generate all variable names
var_names = [f"x_{i}_{j}" for i in range(len(plants)) for j in range(len(markets))]

# Generate variable names with conditions
supply = [100, 150, 200]
large_plants = [plants[i] for i in range(len(plants)) if supply[i] > 120]

# Process solution data
solution_values = [0, 50, 30, 80, 0, 90, 70, 0, 20]
positive_flows = [val for val in solution_values if val > 0]

# More complex comprehensions for optimization
costs = [3, 5, 2, 4, 2, 6, 1, 3, 4]
# Calculate total cost for positive flows only
nonzero_costs = [costs[i] * solution_values[i] 
                 for i in range(len(costs)) 
                 if solution_values[i] > 0]

# Dictionary comprehension for parameter processing
supply_dict = {plants[i]: supply[i] for i in range(len(plants))}

# Set comprehension for unique elements
used_plants = {plants[i//len(markets)] for i, val in enumerate(solution_values) if val > 0}
\end{lstlisting}

\subsection{Functions}
Functions help organize code and create reusable components for optimization workflows.

\begin{lstlisting}[language=Python]
def calculate_transportation_cost(from_plant, to_market, quantity, cost_matrix):
    """Calculate cost for transporting quantity from plant to market."""
    # Simple cost lookup (in practice, you'd have more sophisticated logic)
    unit_cost = cost_matrix.get((from_plant, to_market), 0)
    return unit_cost * quantity

def create_variable_name(plant_id, market_id, prefix="x"):
    """Generate standardized variable names."""
    return f"{prefix}_{plant_id}_{market_id}"

def summarize_solution(variables, values):
    """Create a summary of the optimization solution."""
    solution_dict = {}
    for i, var in enumerate(variables):
        if values[i] > 0:  # Only include positive values
            solution_dict[var] = values[i]
    return solution_dict

def validate_supply_demand_balance(supply, demand):
    """Check if total supply meets total demand."""
    total_supply = sum(supply.values()) if isinstance(supply, dict) else sum(supply)
    total_demand = sum(demand.values()) if isinstance(demand, dict) else sum(demand)
    
    if total_supply >= total_demand:
        return True, "Supply meets demand"
    else:
        shortage = total_demand - total_supply
        return False, f"Shortage of {shortage} units"

# Example usage
cost_matrix = {('Plant1', 'Market1'): 3, ('Plant1', 'Market2'): 5}
transport_cost = calculate_transportation_cost('Plant1', 'Market1', 50, cost_matrix)

var_name = create_variable_name(1, 2)  # Returns "x_1_2"

# Function with default parameters
def format_solution_report(solution_dict, title="Optimization Results"):
    """Format solution dictionary into readable report."""
    report = f"\n{title}\n" + "="*len(title) + "\n"
    for var, value in solution_dict.items():
        report += f"{var}: {value:.2f}\n"
    return report
\end{lstlisting}

\subsection{Essential Package Imports}
Understanding how to import and use key packages is crucial for optimization work.

\begin{lstlisting}[language=Python]
# Standard library imports
import itertools  # For generating combinations, products
from collections import defaultdict  # For grouping data

# Essential numerical computing
import numpy as np
import pandas as pd

# Plotting and visualization
import matplotlib.pyplot as plt

# Geospatial analysis (when working with location-based optimization)
import geopandas as gpd

# Optimization
import pulp

# Example usage of imports
# Generate all plant-market combinations
plants = ['Plant1', 'Plant2']
markets = ['Market1', 'Market2', 'Market3']
all_routes = list(itertools.product(plants, markets))

# Use defaultdict for grouping
solution_by_plant = defaultdict(list)
for route, flow in zip(all_routes, [50, 30, 0, 80, 20, 40]):
    if flow > 0:
        solution_by_plant[route[0]].append((route[1], flow))

# NumPy for numerical operations
flows = np.array([50, 30, 0, 80, 20, 40])
total_flow = np.sum(flows)
nonzero_flows = np.nonzero(flows)[0]

# Pandas for data management
solution_df = pd.DataFrame({
    'From': [route[0] for route in all_routes],
    'To': [route[1] for route in all_routes],
    'Flow': flows
})

active_routes = solution_df[solution_df['Flow'] > 0]
print("Active transportation routes:")
print(active_routes)
\end{lstlisting}

\begin{info}
\textbf{Best Practices for LP Programming:}
\begin{itemize}
\item Use descriptive variable names that reflect the business meaning
\item Organize data in dictionaries or DataFrames for easy access
\item Use list comprehensions to generate sets of similar variables or constraints
\item Create functions for repetitive calculations or data processing
\item Always validate your data (check supply-demand balance, non-negative values, etc.)
\end{itemize}
\end{info}

This foundation prepares you for the optimization modeling work with PuLP that follows. Each concept builds upon the others: you'll use lists and dictionaries to organize problem data, NumPy arrays for numerical computations, pandas DataFrames for data analysis, loops and comprehensions to generate variables and constraints systematically, and functions to organize your optimization workflow.